\section{Specialization and the return to scissors congruence}
\label{cp:sec-specialization}

Theorem~\ref{cp:thm-relative-defect-zero} is a relative identity over the
algebraic parameter space \(U\).  We now explain, without suppressing the
comparison maps, why its fibre at the original tetrahedron proves
Theorem~\ref{cp:thm-main}.

\subsection{The two weight-two injections}

Let \(S_4(\C)\) denote Goncharov's abelian group of oriented smooth
quadrics in \(\mathbf P^3\) equipped with four ordered hyperplane faces,
modulo the relations recalled in Section~\ref{cp:sec-motivic-dictionary}.
Put \(S_j(\C)_\Q=S_j(\C)\otimes_{\mathbf Z}\Q\).  In weight one the
cross-ratio isomorphism identifies
\(S_2(\C)_\Q\) with \(\C^*\otimes_{\mathbf Z}\Q\).  The weight-two part of
the normalized cobar complex begins, in cohomological degrees one and two,
as
\[
 S_4(\C)_\Q\xrightarrow{\ \overline\Delta\ }
 S_2(\C)_\Q\otimes_\Q S_2(\C)_\Q.
\]
There is no degree-zero term.  Consequently its degree-one cohomology is
\begin{equation}
 H^1(S^\bullet(2))_\Q=\ker\overline\Delta.
 \label{cp:eq-primitive-quadric-group}
\end{equation}

We use the following two theorems of Goncharov.

\begin{theorem}[Goncharov's weight-two comparison]
\label{cp:thm-goncharov-weight-two}
There is a canonical map
\begin{equation}
 c_G:H^1(S^\bullet(2))_\Q
 \longrightarrow \gr_2^\gamma K_3(\C)\otimes\Q,
 \label{cp:eq-goncharov-k3-map}
\end{equation}
obtained from the framed relative motive of a quadric simplex.  It is
injective.
\end{theorem}

\begin{proof}
This is the case \(n=2,i=1\) of
\cite[Theorem~4.13, equation~(4.7)]{Goncharov1999}.  The theorem proves in
weight two both the existence and injectivity asserted in Conjecture~4.11
of that paper.
\end{proof}

Let
\[
 \Pred(S^3)=\operatorname{coker}
 \bigl(\Sigma:\Pfull(S^2)\longrightarrow\Pfull(S^3)\bigr)
\]
be the reduced spherical scissors group; \(\Sigma\) is spherical
suspension.  If \(P\) lies in an equatorial \(S^2\subset S^3\), then
\(\Sigma P\) is the double cone obtained by joining \(P\) to the two
antipodal poles.  Cutting and isometries make this construction an additive
map on scissors groups.

\begin{lemma}[Comparison of the two spherical presentations]
\label{cp:lem-spherical-presentations}
Let \(\mathcal P_G(S^3)\) be Goncharov's group generated by oriented,
ordered spherical tetrahedra with his degeneracy, isometry, permutation,
and five-term additivity relations.  There is a canonical isomorphism
\[
 \Pred(S^3)\xrightarrow{\sim}\mathcal P_G(S^3).
\]
\end{lemma}

\begin{proof}
We spell out the comparison because the unreduced and reduced spherical
groups must not be confused.  Let \(C_k(S^3)\) be the free abelian group on
ordered \((k+1)\)-tuples of points of \(S^3\), with boundary
\[
 \partial(x_0,\ldots,x_k)=
 \sum_{i=0}^k(-1)^i(x_0,\ldots,\widehat{x_i},\ldots,x_k).
\]
Let \(D_*\subset C_*(S^3)\) be the subcomplex generated by tuples lying in
a proper great subsphere, and put
\[
 \operatorname{St}(S^3)=H_3\bigl(C_*(S^3)/D_*\bigr).
\]
A proper ordered four-tuple is a cycle in this quotient because each of its
three-point faces lies in a great \(S^2\); a degenerate tuple is zero; and
the boundary of an ordered five-tuple is exactly the alternating five-term
relation.  Recording an ambient orientation and alternating in the four
entries gives Goncharov's permutation rule.  Taking
\(H_0(SO(4),\operatorname{St}(S^3))\), that is, coinvariants \emph{after}
forming this homology group, imposes the oriented-isometry relation.
Allowing \(O(4)\) while transporting the orientation gives the same quotient
by Gerling's theorem.  Thus Goncharov's presentation is
\(H_0(SO(4),\operatorname{St}(S^3))\).

The proof of Dupont's Theorem~7.14 identifies this group canonically with
\[
 \Pfull(S^3)\big/\Sigma\Pfull(S^2)=\Pred(S^3).
\]
The theorem further identifies it with the degree-three homology of the
full coinvariant configuration complex; this further description is not
needed here.  Under the first isomorphism, a proper ordered tuple maps to
its oriented geometric simplex.  This proves the claim
\cite[Theorem~7.14]{Dupont2001}
\cite[Chapter~2, Section~1]{Goncharov1999}.
\end{proof}

\begin{theorem}[Goncharov's spherical map]
\label{cp:thm-goncharov-spherical}
Let
\[
 S_{4,\Q}^{\mathrm{sph}}
 =\left(
    S_4(\C)\otimes_{\mathbf Z}\mathbf Z(2)_{\mathrm{par}}
   \right)^+\otimes_{\mathbf Z}\Q,
 \qquad
 \mathbf Z(2)_{\mathrm{par}}=(2\pi i)^2\mathbf Z,
\]
where complex conjugation acts on both factors and \({}^+\) denotes the
invariant part.  Then complexification of a spherical simplex gives a
canonical injective map
\begin{equation}
 \psi_S(2):\Pred(S^3)\otimes\Q
 \lhook\joinrel\longrightarrow S_{4,\Q}^{\mathrm{sph}}.
 \label{cp:eq-spherical-quadric-map}
\end{equation}
It takes the spherical Dehn differential to the quadric coproduct.
For a tetrahedron with angle Gram matrix \(G\), its value is represented by
the oriented coordinate quadric simplex \(Q_{G^{-1}}\) with its four
coordinate hyperplanes.
\end{theorem}

\begin{proof}
Use Lemma~\ref{cp:lem-spherical-presentations} to identify the source with
Goncharov's group.  The asserted map and its compatibility with the
coproduct are the spherical parts of
\cite[Theorem~4.14]{Goncharov1999}; its injectivity is stated immediately
after the proof of that theorem.  To see the last
description concretely, let \(n_0,\ldots,n_3\) be the four outward facet
normals.  Their Gram matrix is \(G\).  Write a vector in the basis dual to
the \(n_i\), or equivalently use the four facet linear forms
\(v\mapsto\langle n_i,v\rangle\) as coordinates.  In those coordinates the
ambient spherical quadratic form has matrix \(G^{-1}\), and the four
facets are the coordinate hyperplanes.  The determinant root selects the
same family of maximal isotropic planes used to orient Goncharov's
generator.  We choose the ambient ordered orientations of the two real
simplices so that these roots are the common lift
\((G,w)\mapsto(G^R,w)\) of
\eqref{cp:eq-oriented-regge-lift}.  Gerling's theorem permits this
orientation choice without changing the ordinary or reduced spherical
scissors classes.
\end{proof}

The line \(\mathbf Z(2)_{\mathrm{par}}\) only records conjugation parity; it
is not the motivic Tate object \(\Q(2)\) used in
\(H^1(-,\Q(2))\).  After tensoring with \(\Q\) it is one-dimensional and
may be forgotten when testing whether a class vanishes.  Goncharov's
projective-duality map is an involutive linear automorphism of
\(S_4(\C)_\Q\); it commutes with conjugation and hence preserves the
spherical invariant part.  Thus a class vanishes if and only if its dual class,
represented in our coordinates by \(Q_G\), vanishes.

\subsection{Why the relative motive is the coefficient in
\texorpdfstring{\(c_G\)}{cG}}

There is one comparison which must be made literally.  Knowing that some
motivic realization of a quadric class vanishes would not suffice unless it
were the realization occurring in \eqref{cp:eq-goncharov-k3-map}.

\begin{lemma}[Fibre comparison]
\label{cp:lem-fibre-comparison}
Let \(s:\Spec\C\to U\) be a complex point.  Let \(\delta_s^\vee\) be the
primitive difference of the two dual quadric-scissors generators, regarded
as an element of \(\ker\overline\Delta\subset S_4(\C)_\Q\).  The endpoint
class obtained by pulling back the relative framed
defect \(\delta_U^\vee\) is exactly
\[
 c_G(\delta_s^\vee)
 \in H^1(\C,\Q(2))
 =\gr_2^\gamma K_3(\C)\otimes\Q.
\]
\end{lemma}

\begin{proof}
For any smooth base \(T\), write
\[
 \Phi_T(Q,\alpha;M)
 =
 [\,D\widetilde m_T(Q,M)(2);\,v_\alpha,f_M\,]
\]
for the framed coefficient of an oriented quadric simplex.  Here
\(\widetilde m_T(Q,M)\) is the alternating face complex of
Definition~\ref{cp:def-face-motive}; the ruling \(\alpha\) fixes the extreme
frame \(v_\alpha\), and the ordered oriented hyperplanes fix the other
frame \(f_M\).  Interchanging the ruling or applying an odd permutation
changes exactly one frame, so this has Goncharov's sign conventions.

This is not a new realization of his generator.  Chapter~5, Section~2 of
\cite{Goncharov1999} defines
\[
 m(Q,L)=
 \underline{\operatorname{Hom}}\bigl(\widetilde m(Q,L),\Q\bigr)
\]
from the same four-term face complex, with the zero-dimensional faces in
degree \(0\) and \(\mathbf P^3\setminus Q\) in degree \(3\).  Hence no
cohomological shift occurs.  In our convention
\(\Q(1)=\Gm[-1]\), its raw extreme Tate factors are \(\Q(0)\) and
\(\Q(-2)\); tensoring by \(\Q(2)\) rewrites the coefficient as an extension
of \(\Q(0)\) by \(\Q(2)\).  Thus Goncharov's coefficient, in our
normalization, is exactly \(\Phi_T=\mathfrak m_T\).  This is also Brown's
normalization \cite[Definition~5.8]{Brown2013}.  The modern category
\(\DM(T)_\Q\) represents the same bounded complex of smooth strata, its
internal dual, and its two geometric frame maps.

Every face term is constructible and dualizable.  Exact symmetric monoidal
pullback therefore gives a canonical framed isomorphism
\begin{equation}
 s^*\Phi_U(Q,\alpha;M)
 \simeq
 \Phi_\C(Q_s,\alpha_s;M_s).
 \label{cp:eq-pullback-coefficient}
\end{equation}
Indeed, pullback commutes term by term with the incidence maps, internal
duality, the Tate twist, and the frames.  It also preserves the explicit
face filtration and its endpoint connecting morphism.  No motivic
\(t\)-structure over \(\C\), and no \(t\)-exactness assertion, is needed.

A primitive coefficient determines an endpoint triangle
\begin{equation}
 \Q_U(2)\longrightarrow E\longrightarrow\Q_U(0)
 \xrightarrow{\ \kappa_U\ }\Q_U(2)[1].
 \label{cp:eq-endpoint-triangle}
\end{equation}
Equations~(4.6)--(4.7) of \cite{Goncharov1999} first send a
quadric-simplex generator to \(\Phi_\C\), then apply degree-one normalized
cobar cohomology.  In weight two the resulting map is
\begin{equation}
\begin{split}
 H^1(\Phi_\C):H^1(S^\bullet(2))_\Q
 &\longrightarrow
 \operatorname{Hom}_{\DM(\C)_\Q}
       \bigl(\Q(0),\Q(2)[1]\bigr)\\
 &\cong\gr_2^\gamma K_3(\C)\otimes\Q,
\end{split}
\label{cp:eq-coefficient-cobar-map}
\end{equation}
and this is the map \(c_G\) of
Theorem~\ref{cp:thm-goncharov-weight-two}.  Because the quadric cobar
complex starts in degree \(1\), a primitive element of \(S_4(\C)_\Q\) is
literally a degree-one cocycle.  On such a cocycle
\eqref{cp:eq-coefficient-cobar-map} is its endpoint connecting morphism:
the reduced coproduct is exactly the obstruction contributed by the middle
Tate degree.
Here \(H^1(\Phi_\C)\) is shorthand for this direct operation on primitive
degree-one cocycles; we are not asserting a global abelian mixed-Tate heart
over \(\C\).

Apply this construction to the additive difference
\(\delta_U^\vee\).  Equation~\eqref{cp:eq-pullback-coefficient} identifies
its fibre with the coefficient of \(\delta_s^\vee\), and naturality of the
endpoint triangle gives
\[
 s^*\kappa_U(\delta_U^\vee)
 =\kappa_\C\bigl(\Phi_\C(\delta_s^\vee)\bigr)
 =H^1(\Phi_\C)(\delta_s^\vee)
 =c_G(\delta_s^\vee).
\]
This is an equality in motivic cohomology before any regulator is applied.

Finally, two different dualities occur in the notation.  The superscript
\(\vee\) is projective duality, the antipode taking
\(Q_{G^{-1}}\) to \(Q_G\) and transporting its ruling.  The \(D\) inside
\(\Phi_T\) is the internal cohomological dual of the resulting face
complex.  Thus no second unrecorded projective dual or sign is present.
Theorem~4.13 proves the existence and injectivity of
\eqref{cp:eq-coefficient-cobar-map} in weight two.  Goncharov's stronger
Conjecture~4.15 asks for an explicit chain map from the entire quadric
complex to a Bloch complex; neither that chain map nor any Bloch-symbol
formula is used here.
\end{proof}

\subsection{Specializing to the physical tetrahedron}

Let \(T,T^R\) be the pair in Theorem~\ref{cp:thm-main}.  By
Proposition~\ref{cp:prop-physical-point-on-chart}, its physical angle phases
and positive determinant root define a point
\[
 s_T:\Spec\C\longrightarrow U
\]
on one of the two rational charts.  Put
\[
 z=[T]-[T^R]\in\Pfull(S^3),
 \qquad \bar z\in\Pred(S^3).
\]

The spherical map sends
\(\bar z\otimes1\) to
\(\delta_{s_T}\otimes(2\pi i)^2\).  Forgetting this nonzero
one-dimensional parity factor, we regard \(\delta_{s_T}\) as an element of
\(S_4(\C)_\Q\).  Projective duality sends it to
\(\delta_{s_T}^\vee\), the fibre of the relative defect considered in
Section~\ref{cp:sec-relative-rigidity}.  Lemma
\ref{cp:lem-coproduct-zero}, in its quadric-scissors form, makes it
primitive.

Theorem~\ref{cp:thm-relative-defect-zero} says that
\(\delta_U^\vee=0\) as a framed class.  Therefore its endpoint morphism is
zero, and Lemma~\ref{cp:lem-fibre-comparison} gives
\[
 c_G(\delta_{s_T}^\vee)=0.
\]
The injectivity of \(c_G\) in Theorem
\ref{cp:thm-goncharov-weight-two} implies
\(\delta_{s_T}^\vee=0\) in \(S_4(\C)_\Q\), and projective duality then gives
\(\delta_{s_T}=0\) there.  Injectivity of \(\psi_S(2)\) yields
\begin{equation}
 \bar z=0\qquad\text{in }\Pred(S^3)\otimes\Q.
 \label{cp:eq-reduced-regge-zero-rational}
\end{equation}

This completes the motivic part of the proof.  The remaining paragraph is
ordinary spherical scissors geometry.

\subsection{Removing rationalization and suspension}

We record the three standard spherical scissors facts used here
\cite[Example~7.8 and Corollary~9.19(b)]{Dupont2001}:
\begin{enumerate}[label=\textup{(S\arabic*)},leftmargin=2.7em]
\item \(\Pfull(S^3)\) is uniquely divisible, hence is a \(\Q\)-vector
space.
\item Area is an isomorphism
\[
 \Area:\Pfull(S^2)\xrightarrow{\sim}\R.
\]
\item Equality in the spherical scissors group can be converted into a
finite equidecomposition by Zylev cancellation; Gerling's theorem permits
the motions to be taken in \(SO(4)\)
\cite[after Definition~1.6 and Theorem~2.2]{Dupont2001}.
\end{enumerate}

Since \(\Pfull(S^2)\cong\R\) and \(\Pfull(S^3)\) are \(\Q\)-vector
spaces, the image of suspension is a \(\Q\)-subspace.  Its cokernel
\(\Pred(S^3)\) is therefore also a \(\Q\)-vector space.  Consequently
\eqref{cp:eq-reduced-regge-zero-rational} already says
\(\bar z=0\) before tensoring with \(\Q\).  By the definition of the
reduced group, there is a class \(Q\in\Pfull(S^2)\) with
\begin{equation}
 z=\Sigma Q.
 \label{cp:eq-regge-suspension}
\end{equation}

If a spherical polygon of area \(A\) is suspended from the two poles of
\(S^3\), the warped-product volume element is
\(\cos^2(t)\,dt\,dA\), where \(-\pi/2\leq t\leq\pi/2\).  Therefore
\begin{equation}
 \Vol(\Sigma Q)
 =\left(\int_{-\pi/2}^{\pi/2}\cos^2t\,dt\right)\Area(Q)
 =\frac\pi2\Area(Q).
 \label{cp:eq-suspension-volume}
\end{equation}
Both sides are additive under cutting, so
\eqref{cp:eq-suspension-volume} holds for every class
\(Q\in\Pfull(S^2)\), not only for a single polygon.
Regge mates have equal volume by Theorem~\ref{cp:thm-regge-input}.  Applying
volume to \eqref{cp:eq-regge-suspension} gives \(\Area(Q)=0\).  The area
classification forces \(Q=0\), so \(z=0\) in \(\Pfull(S^3)\).

Zylev cancellation and Gerling's theorem now turn this equality into a
finite equidecomposition.  This proves Theorem~\ref{cp:thm-main}.
